\section{Actual tangent coordinates are regular symbols}

Write $S=\mathbb C[[z_1,\ldots,z_n]]$, $\mathfrak m=(z_1,\ldots,z_n)$,
and $\operatorname{in}(K)$ for the entire homogeneous initial ideal of
an ideal $K\subset S$. All powers below are powers of the actual maximal
ideal; none is a normal-power filtration along a positive-dimensional
centre. The conclusions supply the finite-order regular-symbol step used
in comparing those filtrations, without asserting normal flatness.

\begin{theorem}[Initial tangent stability gives actual formal regular symbols]
\label{hir:actual-initial-tangent-formal-regular-symbol}
\uses{hir:actual-polynomial-pderiv-locally-nilpotent,hir:actual-pderiv-stable-polynomial-coordinate-saturation,hir:actual-homogeneous-initial-ideal-degree-lift,hir:actual-homogeneous-initial-component-derivative-dictionaries}
\lean{Hironaka.HironakaActualInitialSymbolRegularity.formal_regular_symbol_of_initial_pderiv}
\leanok
Suppose the polynomial contraction of $\operatorname{in}(K)$ is stable
under $\partial/\partial z_i$. For every $q\geq0$ and every
$b\in\mathfrak m^q$,
\[
 z_i b\in K+\mathfrak m^{q+2}\quad\Longrightarrow\quad
 b\in K+\mathfrak m^{q+1}.
\]
\end{theorem}
\begin{proof}
Write $z_i b=g+c$, with $g\in K$ and
$c\in\mathfrak m^{q+2}$. Both $z_i b$ and $c$ belong to
$\mathfrak m^{q+1}$, so $g\in\mathfrak m^{q+1}$. Its degree
$q+1$ component therefore belongs to $\operatorname{in}(K)$.
The same component of $c$ is zero, and multiplication by $z_i$
raises degree by one. Consequently $z_i b_q\in\operatorname{in}(K)$,
where $b_q$ denotes the degree $q$ homogeneous component of $b$.
The component $b_q$ is an actual polynomial. Partial stability of the
polynomial contraction and characteristic-zero coordinate saturation give
$b_q\in\operatorname{in}(K)$.

The homogeneous initial-ideal degree-lifting theorem supplies
$L\in K\cap\mathfrak m^q$ whose degree $q$ component is $b_q$.
Every coefficient of $b-L$ of degree below $q$ vanishes because both
summands belong to $\mathfrak m^q$; its coefficients of degree $q$
vanish by the component equality. Thus $b-L\in\mathfrak m^{q+1}$,
and $b\in K+\mathfrak m^{q+1}$ as required.
\end{proof}

\begin{theorem}[The same symbol is regular in the original holomorphic ring]
\label{hir:actual-initial-tangent-original-regular-symbol}
\uses{hir:actual-initial-tangent-formal-regular-symbol,hir:actual-formal-extension-whole-original-ideal-contraction,hir:actual-original-formal-variable-power-extension}
\lean{Hironaka.HironakaActualInitialSymbolRegularity.original_regular_symbol_of_initial_pderiv,Hironaka.HironakaActualInitialSymbolRegularity.unshifted_original_regular_symbol_of_initial_pderiv}
\leanok
Let $I$ be an ideal in the original ring of holomorphic germs at the
origin, and let $K$ be its entire Taylor extension. If the polynomial
contraction of $\operatorname{in}(K)$ is stable under the corresponding
coordinate partial derivative, multiplication by the actual coordinate
germ $v=z_i$ has the preceding regular-symbol property for every degree
of the original maximal-ideal filtration of $\mathcal O/I$.
\end{theorem}
\begin{proof}
The actual Taylor algebra homomorphism sends the original coordinate
germ to $z_i$, and sends every power of the original maximal ideal
onto the corresponding formal variable-ideal power. It therefore sends
$vb\in I+\mathfrak m^{q+2}$ and $b\in\mathfrak m^q$ to the exact
hypotheses of the formal theorem. That theorem puts the Taylor image of
$b$ in the entire extension of $I+\mathfrak m^{q+1}$.
Contraction of the entire original ideal under the Taylor homomorphism
gives $b\in I+\mathfrak m^{q+1}$. Reindexing the identical dimension
$0+n=n$ preserves both the coordinate and ideal dictionaries, so the
statement also holds in the original unshifted coordinate ring.
\end{proof}

\begin{theorem}[Actual tangent coordinates are nonzerodivisors on the original quotient]
\label{hir:actual-initial-tangent-original-coordinate-saturation}
\uses{hir:actual-initial-tangent-original-regular-symbol,hir:regular-symbol-whole-successor-division,hir:actual-whole-ideal-maximal-filtration-separated}
\lean{Hironaka.HironakaActualInitialSymbolRegularity.original_coordinate_mul_mem_iff_of_initial_pderiv}
\leanok
Under the preceding hypotheses, every actual original germ $a$ satisfies
$va\in I$ if and only if $a\in I$. Thus multiplication by the actual
coordinate class is injective on the entire original quotient
$\mathcal O/I$.
\end{theorem}
\begin{proof}
If $a\in I$, ideal closure gives $va\in I$. Conversely suppose
$va\in I$. For every $q\geq0$ this implies
$va\in I+\mathfrak m^{q+2}$. Strict successor division gives
$a\in I+\mathfrak m^{q+1}$. The original analytic germ ring is
Noetherian and local, so the maximal-adic filtration of its quotient
by $I$ is separated. Intersecting these positive thickenings gives
$a\in I$. Applying this equivalence to the difference of two
representatives proves the asserted injectivity on the entire
quotient.
\end{proof}

\begin{theorem}[Original Hilbert--Samuel constant centres give strict tangent division]
\label{hir:actual-original-hs-centre-tangent-strict-division}
\uses{hir:actual-original-hs-centre-whole-initial-tangent-pderiv,hir:actual-initial-tangent-original-regular-symbol,hir:actual-initial-tangent-original-coordinate-saturation,hir:regular-symbol-whole-successor-division,hir:regular-symbol-whole-principal-intersection}
\lean{Hironaka.HironakaActualInitialSymbolRegularity.exists_original_hs_constant_centre_strict_tangent_division}
\leanok
Let $A$ be a coherent ideal on a complex manifold, and let $C$ be
regular of dimension $d$ at $x$. Suppose the entire intrinsic
Hilbert--Samuel function of $A$ is constant along $C$. There is an
actual local holomorphic coordinate change flattening $C$ at $x$
such that the following holds in the unchanged origin ideal $I_0$
of the flattened original sheaf. For each of its first $d$ tangent
coordinate germs $v$, every $q\geq0$, and every original germ $a$,
\[
 va\in I_0+\mathfrak m^{q+1}
 \quad\Longleftrightarrow\quad a\in I_0+\mathfrak m^q,
 \qquad
 (I_0+\mathfrak m^{q+1})\cap(I_0+(v))
   =I_0+(v)\mathfrak m^q.
\]
Moreover $va\in I_0$ if and only if $a\in I_0$, so each such actual
tangent coordinate acts injectively on the entire original quotient.
\end{theorem}
\begin{proof}
The actual original constant-Hilbert--Samuel centre theorem constructs
one flattening and proves stability of the polynomial contraction of
the entire initial ideal under every tangent partial derivative.
Apply the preceding original-ring theorem to each actual tangent
coordinate. It gives the regular-symbol implication in every degree
for the same $I_0$. Each tangent coordinate belongs to the actual
maximal ideal, as follows either by evaluation at the origin or by
the exact Taylor ideal dictionary. Strict filtered successor division
then gives the displayed equivalence for arbitrary original
numerators. Strict principal-image intersection gives the displayed
equality of entire ideals. These are finite-power identities; no
assertion of normal flatness is needed. The original-ring coordinate
saturation theorem gives the last assertion from the proved
finite-power statements and actual analytic separatedness.
\end{proof}
